\chapter{Dataset, Experimental Context, and Data Audit}
\label{ch:dataset}

This chapter sets out the physical experimental context, the structure of
the available data, and the outcome of a dataset-level audit that later
chapters build upon rather than re-deriving. It distinguishes between facts
about the physical experiment, which originate from predecessor work
external to the present project, and facts about how that experimental
record has been organised, cleaned, and represented within the present
project's own pipeline.

\section{Experimental Context}
\label{sec:experimental-context}

The dataset originates from a physical durability study of ferrocement
specimens subjected to accelerated chloride-induced corrosion. Ferrocement
elements reinforced with steel mesh were fabricated, exposed to a saline
environment for an extended period, and photographed repeatedly to record
the progression of visible surface corrosion. At the end of each specimen's
exposure period, a terminal mechanical test was carried out to characterise
its residual structural condition. The physical fabrication, exposure, and
testing programme was designed and executed as predecessor work, documented
in an earlier MSc thesis \citep{hossain_thesis} and a related conference
paper \citep{artiste2025}; the present project did not conduct this physical
experiment. It instead works from the photographic record, the associated
specimen/exposure metadata, and the terminal mechanical-test measurements
that resulted from it, and reconstructs a specimen-to-campaign/treatment
mapping from specimen identifiers, cross-checked against the treatment
tables reported in the predecessor thesis.\footnote{The specimen-to-campaign
and specimen-to-treatment mapping used throughout the present project's
pipeline is maintained as a single, version-controlled table that is
documented as the preferred source for this reconstruction, in preference
to the raw treatment columns present in the data workbook.}
Facts concerning the physical experiment itself are attributed to this
predecessor source; facts concerning how the present project has organised
or audited the resulting data are attributed to the present project's own
code, configuration, or generated outputs.

The experimental programme comprises two campaigns, executed at different
times and, notably, at different steel-mesh and chloride-exposure settings.
The first campaign involved 24 specimens organised into five treatment
series, reinforced with a higher steel-mesh configuration and exposed to a
lower chloride concentration; specimens from this campaign were
terminally tested after 196 days of ageing (week 28 of observation). The
second campaign involved a further 24 specimens organised into four
treatment series, reinforced with a lower steel-mesh configuration and
exposed to a higher chloride concentration; specimens from this campaign
were terminally tested after 252 days of ageing (week 36 of observation).
Table~\ref{tab:campaign-design} summarises the resulting design matrix, and
Section~\ref{sec:confounding} discusses a structural consequence of this
design that later chapters must account for when interpreting results.

Across both campaigns, seven coarse treatment protocols are represented: an
untreated control; a corrosion inhibitor mixed into the mortar matrix at
casting; a corrosion inhibitor applied to the surface after casting
(applied by different methods — general surface application in the first
campaign, spray or brush application in different series of the second
campaign); a protective surface paint; a combined paint-plus-surface-inhibitor
treatment; and two treatments incorporating a glass-based surface compound,
applied either alone or together with the surface inhibitor. These seven
protocols are unevenly represented, both overall and within each campaign,
as shown in Table~\ref{tab:campaign-design}.

Terminal structural condition was assessed through a single mechanical test
per specimen. Two structural quantities are reported in the predecessor
documentation: the peak load recorded during a four-point bending test to
the specimen's load-bearing limit, and the loss of cross-sectional area of
the outermost longitudinal reinforcing wires near the aged surface,
obtained from microscopic examination of a cross-section taken at the
failure location \citep{hossain_thesis}.\footnote{This attribution follows
the present project's own audit documentation, which cites the relevant
thesis chapters for these two definitions; the thesis itself was not
re-read in full for the present chapter.}
Both quantities are recorded once per specimen, at the single terminal time
point associated with that specimen's campaign.

\begin{table}[htbp]
\centering
\caption{Experimental campaign and design matrix. Steel-mesh configuration,
chloride concentration, and terminal testing time are recorded once per
specimen and do not vary within a campaign. Treatment-group specimen
counts are reconstructed from the project's specimen-mapping table.}
\label{tab:campaign-design}
\begin{tabular}{lccccc}
\toprule
Campaign & Specimens & Steel-mesh & NaCl & Terminal week & Terminal ageing \\
 & & configuration & concentration & (observation) & (days) \\
\midrule
Campaign 1 & 24 & 7 & 3.5\% & 28 & 196 \\
Campaign 2 & 24 & 4 & 5.0\% & 36 & 252 \\
\bottomrule
\end{tabular}

\vspace{2mm}
\footnotesize
\setlength{\tabcolsep}{3.5pt}
\begin{tabular}{lcccccccc}
\toprule
& \multicolumn{7}{c}{Specimen count by treatment protocol} & \\
\cmidrule(lr){2-8}
Campaign & Control & Mixed-in & Surface & Paint & Paint+ & Surface+ & Glass & Total \\
 & (NO) & inhibitor & inhibitor & (PA) & surface & glass & compound & \\
 & & (MI) & (SA) & & (SA\_PA) & (SA\_VF) & (VF) & \\
\midrule
Campaign 1 & 4 & 6 & 6 & 1 & 3 & 2 & 2 & 24 \\
Campaign 2 & 6 & -- & 12 & 6 & -- & -- & -- & 24 \\
\midrule
Total & 10 & 6 & 18 & 7 & 3 & 2 & 2 & 48 \\
\bottomrule
\end{tabular}
\normalsize
\end{table}

\section{Data Sources and Modalities}
\label{sec:data-modalities}

The material available to the present project is organised into four
distinct modalities, which differ substantially in their granularity and
availability and must not be conflated.

\textbf{A. Image observations.} Photographs of each specimen's surface,
taken repeatedly over its observation period. Each image is associated
with a specimen identifier, an acquisition date, and an observation week,
encoded in the image filename. This is the only modality that is
genuinely longitudinal at fine time resolution (Section~\ref{sec:longitudinal}).

\textbf{B. Specimen and exposure metadata.} Design and exposure variables
recorded once per specimen and constant across that specimen's entire
observation period: campaign identity, treatment series and protocol,
steel-mesh configuration, and chloride concentration (Table~\ref{tab:campaign-design}).
A small number of additional fields, such as observation week itself, vary
by image rather than by specimen.

\textbf{C. Visible-corrosion labels.} Manually assigned quantitative and
categorical measures of surface corrosion, recorded for essentially every
usable image: a continuous total (whole-specimen) rust percentage, a
continuous peak (localised maximum) rust percentage, and, in the currently
audited data table, a four-class ordinal severity category for each of
these two measures. These labels are therefore available at the same fine
time resolution as the images themselves. As discussed in
Section~\ref{sec:dataset-audit}, the number of ordinal levels used for
these severity categories is itself not constant across the project's
successive dataset builds.

\textbf{D. Structural endpoint measurements.} The two terminal mechanical
quantities introduced in Section~\ref{sec:experimental-context} — wire-area
loss and ultimate load — recorded once per specimen, at that specimen's
single terminal observation.

The essential asymmetry that this chapter establishes, and that
Section~\ref{sec:structural-supervision} develops in detail, is that
modalities A and C are dense (available for nearly every image
observation), while modality D is sparse (available for only one
observation per specimen). Modality B sits between the two: it is
constant per specimen but known for every observation of that specimen.
Table~\ref{tab:data-modalities} summarises the granularity and
availability of each modality, including the class distributions of the
two four-class visible-corrosion severity categories, which are markedly
imbalanced toward the lowest-severity class.

\begin{table}[htbp]
\centering
\caption{Data modalities and label availability, computed directly from
the aligned 791-row dataset table.}
\label{tab:data-modalities}
\begin{tabular}{p{2.6cm}p{5.3cm}p{2.0cm}p{3.4cm}}
\toprule
Modality & Content & Granularity & Availability \\
\midrule
A. Image observations & Specimen surface photograph & Per image & 791 of
792 aligned observations \\
B. Specimen/exposure metadata & Campaign, treatment, steel-mesh
configuration, chloride concentration, observation week & Per specimen
(design variables); per image (week) & All 791 aligned observations \\
C. Visible-corrosion labels & Total and peak rust percentage (continuous);
four-class severity category for each & Per image & 791 of 792 aligned
observations. Total-category counts: 658 / 99 / 20 / 14 (classes 1--4);
peak-category counts: 526 / 107 / 95 / 63 \\
D. Structural endpoint measurements & Wire-area loss; ultimate load & Per
specimen (single terminal observation) & 48 of 791 aligned observations
(one per specimen) \\
\bottomrule
\end{tabular}
\end{table}

\section{Longitudinal Observation Structure}
\label{sec:longitudinal}

Each specimen was photographed repeatedly over its observation period,
yielding, after the alignment audit described in
Section~\ref{sec:dataset-audit}, between 15 and 18 image observations per
specimen (median 16). Observation weeks are not sampled on a strict fixed
weekly cadence: across the aligned dataset, images were captured at
twenty-five distinct week values spanning week 0 to week 36, with
irregular spacing between successive observations of the same specimen.

The repeated images belonging to a given specimen are not independent
observations in the sense required by a naive statistical treatment. They
share a single physical specimen, a single fabrication and treatment
history, and a single acquisition and imaging context, and they trace out
one continuous corrosion trajectory rather than twenty-five unrelated
draws from a population of unrelated images. An observation at week 12 of
a given specimen is informative about that same specimen's observation at
week 15 in a way that it is not informative about a different specimen's
observation at week 15.

This has a direct implication for how the data may legitimately be used
to evaluate a predictive model, which the present chapter states only in
general terms: any evaluation procedure that treats individual image rows
as exchangeable, independent units risks allowing information about a
specific specimen to pass from a training partition to a test partition
through a different image of the same specimen, rather than through
genuine generalisation to unseen specimens. The specific evaluation
procedures the project adopted in response to this concern are the subject
of Chapter~\ref{ch:framework}.

\section{Structural Supervision}
\label{sec:structural-supervision}

The dataset comprises 48 specimens in total. Each specimen contributes
exactly one structural observation — its terminal wire-area-loss and
ultimate-load measurement, taken at that specimen's single terminal test
(week 28 for Campaign~1 specimens, week 36 for Campaign~2 specimens; see
Table~\ref{tab:campaign-design}). No specimen contributes more than one
structural observation, and no specimen's structural condition is measured
at any intermediate observation week.

This makes structural supervision fundamentally different in kind from the
dense visible-corrosion labels described in
Section~\ref{sec:data-modalities}. Visible-corrosion labels describe a
trajectory: for a given specimen, they are available at many points along
its observation period, and changes in visible corrosion over time can be
observed directly. Structural measurements describe a single endpoint:
for a given specimen, nothing is known about its structural condition
before the terminal observation, and nothing is known about it after.

A direct consequence is that the dataset does not contain the information
required to pose intermediate structural condition, a structural
trajectory, or a time-to-failure quantity as a directly supervised
learning problem. Doing so would require either repeated structural
measurements of the same specimen over time, or a recorded failure event
time, and neither is present. Table~\ref{tab:quantity-availability}
summarises this distinction, separating quantities that are directly
observed in the data from quantities that later phases of the project
construct through modelling, quantities that are further derived from
those modelled quantities for screening purposes, and quantities that are
not available at all.

\begin{table}[htbp]
\centering
\caption{Directly observed, model-derived, derived-screening, and
unavailable quantities. This distinction is maintained throughout the
report and must not be blurred: a model-derived estimate of intermediate
structural condition is not a measurement, and a derived screening output
is not a supervised prediction target.}
\label{tab:quantity-availability}
\begin{tabular}{p{3.1cm}p{9.6cm}}
\toprule
Category & Quantities \\
\midrule
Directly observed &
Image; observation week; visible-corrosion labels (total and peak rust
percentage and their four-class categories); terminal structural
measurements (wire-area loss, ultimate load) at each specimen's single
terminal observation. \\[2mm]
Model-derived &
Estimated hidden structural condition at intermediate (non-terminal)
observation weeks; fitted specimen-level degradation trajectories over
the observation period. \\[2mm]
Derived screening outputs &
Threshold-crossing time computed from a fitted degradation trajectory
(``proxy remaining-useful-life''); see Section~\ref{sec:proxy-rul}. \\[2mm]
Not available &
A directly supervised remaining-useful-life or time-to-failure target;
an experimentally observed failure event time. No evidence encountered in
the reviewed project documentation or the predecessor thesis material
establishes such a target for this dataset. \\
\bottomrule
\end{tabular}
\end{table}

\begin{figure}[htbp]
\centering
\resizebox{0.7\textwidth}{!}{% Report schematic (not an experimental result / not to scale).
% Illustrates the asymmetry between dense image-level visible-corrosion
% supervision and sparse specimen-level structural supervision.
% Counts shown are the aligned-dataset figures established in the dataset-audit section.
\begin{tikzpicture}[
    node distance=8mm and 14mm,
    box/.style={draw, rounded corners, align=center, minimum width=5.2cm,
                minimum height=1.1cm, font=\small, fill=gray!6},
    root/.style={draw, rounded corners, align=center, minimum width=4.2cm,
                minimum height=1.1cm, font=\small\bfseries, fill=gray!15},
    arr/.style={-{Latex[length=2.2mm]}, thick},
    lbl/.style={font=\scriptsize\itshape, align=center, text width=5.2cm}
]

\node[root] (root) {48 specimens\\(2 experimental campaigns)};

\node[box, below left=of root, xshift=-1.7cm] (imgs) {repeated photographic\\observations per specimen\\(15--18 images each, weeks 0--36)};
\node[box, below=of imgs] (corr) {visible-corrosion labels\\available for 791 of 792\\image observations};
\node[lbl, below=2mm of corr] (corrlbl) {dense, image-level supervision};

\node[box, below right=of root, xshift=1.7cm] (test) {terminal mechanical test\\(one four-point bending test\\per specimen, week 28 or 36)};
\node[box, below=of test] (struct) {structural measurements\\available for 48 of 791\\rows (one per specimen)};
\node[lbl, below=2mm of struct] (structlbl) {sparse, specimen-level supervision};

\draw[arr] (root.south) -| (imgs.north);
\draw[arr] (imgs) -- (corr);
\draw[arr] (root.south) -| (test.north);
\draw[arr] (test) -- (struct);

\end{tikzpicture}
}
\caption{Schematic contrast between dense, image-level visible-corrosion
supervision and sparse, specimen-level structural supervision, following
directly from Table~\ref{tab:quantity-availability}. This is a report
schematic illustrating the data structure described in
Sections~\ref{sec:data-modalities}--\ref{sec:structural-supervision}; it is
not a plot of experimental results.}
\label{fig:dataset-supervision-schematic}
\end{figure}

\section{Dataset Audit and Alignment}
\label{sec:dataset-audit}

The image archive contains 792 image files in total, spanning 48
specimens. Of these, one file — a single image belonging to a Campaign~2
specimen from the surface-inhibitor (spray) series — is corrupted and
cannot be decoded by standard image-reading routines. The remaining 791
files are readable and align to a corresponding metadata record.
Table~\ref{tab:dataset-overview} summarises the resulting dataset
dimensions on the aligned (791-row) basis that this report adopts by
default.

\begin{table}[htbp]
\centering
\caption{Dataset overview, aligned (791-row) basis, computed directly from
the current canonical dataset table.}
\label{tab:dataset-overview}
\begin{tabular}{ll}
\toprule
Quantity & Value \\
\midrule
Image files (raw archive) & 792 \\
Corrupted / unreadable image files & 1 \\
Aligned rows adopted by this report & 791 \\
Specimens & 48 \\
Experimental campaigns & 2 \\
Observation week range & 0--36 (irregular; 25 distinct week values) \\
Image observations per specimen & 15--18 (median 16) \\
Specimens with structural endpoint measurements & 48 of 48 (one
observation each) \\
\bottomrule
\end{tabular}
\end{table}

The project's own row-count convention for this one corrupted file has
evolved across its successive implementation phases, and this evolution
is best understood as a progressive tightening of data-quality policy
rather than as an unresolved inconsistency. Table~\ref{tab:corruption-handling}
summarises the three distinct policies adopted. The earliest exploratory
implementation removed the affected row from its working table without
substitution, arriving at 791 usable rows. The next three implementation
phases each instead retained all 792 rows and substituted the missing
image-derived quantities — handcrafted features in one phase, a learned
embedding in another, a broader interpretable feature set in a third —
with the corresponding column-wise median value, so that the corrupted
observation remained present but carried an imputed rather than a
measured image representation. The current and subsequent implementation
phases instead apply an explicit, configuration-level data-quality rule
that excludes any image that cannot be aligned to a readable, matched
record, which is the policy this report adopts by default: 791 aligned
rows drawn from 792 image files, 48 specimens.

\begin{table}[htbp]
\centering
\caption{Handling of the single corrupted image observation across
successive implementation phases, verified directly against each phase's
own data-loading code.}
\label{tab:corruption-handling}
\begin{tabular}{p{3.4cm}p{4.0cm}p{5.3cm}}
\toprule
Phase & Resulting row count & Handling mechanism \\
\midrule
Earliest exploratory implementation & 791 & Row silently omitted during
table construction; no substitution. \\
Classical baseline & 792 & Row retained; affected handcrafted image
features replaced by column-wise medians. \\
Deep-embedding phase & 792 & Row retained; affected learned image
embedding replaced by column-wise medians. \\
Interpretable multi-family feature phase & 792 & Row retained and flagged;
affected feature set (91 features) replaced by column-wise medians. \\
Current canonical phase (adopted here), inherited by the later
classification/augmentation effort & 791 & Row excluded by an explicit
alignment/readability policy prior to feature extraction. \\
\bottomrule
\end{tabular}
\end{table}

Because this row-count convention differs across phases, any quantitative
result drawn from an earlier phase and discussed later in this report
(Chapter~\ref{ch:activities}) is understood to be computed on a 792-row
basis with one imputed observation, whereas results drawn from the current
canonical phase and the later classification effort are computed on a
791-row aligned basis. The two bases are not row-for-row identical and are
not treated as interchangeable when results are compared
(Chapter~\ref{ch:results}).

A second, distinct schema change concerns the visible-corrosion severity
category labels introduced in Section~\ref{sec:data-modalities}. The
currently audited data table represents both the total-rust and peak-rust
severity categories on a four-level ordinal scale. An earlier canonical
dataset build, produced during the interpretable multi-family feature
phase on the historical 792-row basis, instead represented the same two
quantities on a five-level ordinal scale, with materially different class
counts (Table~\ref{tab:category-schema}). The data-loading code from that
earlier phase confirms the change directly: it is written to match a raw
workbook column defined over a five-level range, where the currently
audited workbook instead defines the corresponding column over a
four-level range. This is a genuine change in the underlying label schema
between the workbook snapshot used at that phase and the workbook snapshot
audited here, not a data-quality defect internal to the project, and the
two schemas are not treated as interchangeable in this report: a
categorical-classification result reported for the earlier phase in
Chapter~\ref{ch:activities} refers to a five-class problem and is not
described as an evaluation of the four-class scheme that the current
dataset, and the later classification and augmentation effort
(Section~\ref{sec:representation}, Chapter~\ref{ch:activities}), target.

\begin{table}[htbp]
\centering
\caption{Visible-corrosion severity category schema across dataset
builds. Class counts are ordered from the lowest to the highest severity
level.}
\label{tab:category-schema}
\begin{tabular}{p{4.2cm}p{2.1cm}p{3.1cm}p{3.1cm}}
\toprule
Dataset build & Row-count basis & Total-rust category counts & Peak-rust
category counts \\
\midrule
Interpretable multi-family feature phase (historical) & 792 (imputed) &
5 classes: 493 / 51 / 38 / 77 / 133 & 5 classes: 304 / 159 / 72 / 122 /
135 \\
Current canonical phase (audited; adopted by this report) & 791
(aligned) & 4 classes: 658 / 99 / 20 / 14 & 4 classes: 526 / 107 / 95 /
63 \\
\bottomrule
\end{tabular}
\end{table}

\section{Dataset Structure and Potential Confounding}
\label{sec:confounding}

Table~\ref{tab:campaign-design} shows that steel-mesh configuration,
chloride concentration, and terminal ageing duration do not vary
independently of campaign identity within the observed sample: every
specimen with the higher steel-mesh configuration also has the lower
chloride concentration and the shorter terminal ageing duration, and
belongs to Campaign~1; every specimen with the lower steel-mesh
configuration also has the higher chloride concentration and the longer
terminal ageing duration, and belongs to Campaign~2. No specimen in the
dataset combines a Campaign-1 mesh configuration with a Campaign-2
chloride level, or vice versa.

This structural property of the experimental design has a direct
consequence for how any statistical association between a design-linked
quantity and a measured outcome must be interpreted. Where an association
is estimated by pooling specimens across both campaigns, that association
cannot, on the evidence of the pooled sample alone, be uniquely attributed
to any one of the co-varying factors — campaign identity, mesh
configuration, chloride exposure, and ageing duration are, within this
sample, different labels for largely the same partition of specimens. A
pooled correlation involving one of these factors, or a quantity correlated
with one of them, should accordingly be read as a description of the
pooled sample rather than as evidence isolating the effect of that factor.
This report returns to the practical consequences of this property for
model evaluation in Chapters~\ref{ch:framework} and~\ref{ch:results}; the
present section deliberately stops short of discussing any modelling
result.

A related consequence concerns what it means to evaluate a model by
holding out an entire campaign. Because campaign identity coincides with a
specific combination of mesh configuration, chloride level, and ageing
duration, holding out one campaign does not test a model against unseen
specimens drawn from the same underlying design space as its training
data; it tests the model against specimens whose design characteristics
differ systematically, and simultaneously along several correlated
dimensions, from every specimen the model was trained on. Such an
evaluation is better characterised as a test of extrapolation across the
experimental design space than as an ordinary test of generalisation to
unseen specimens within that space. This distinction is developed formally
in Section~\ref{sec:eval-regimes}.

\section{Summary of Dataset Constraints}
\label{sec:dataset-summary}

The dataset available to this project is constrained in ways that
directly shape the modelling strategies described in later chapters:

\begin{itemize}
\item Visible-corrosion supervision is dense (available for 791 of 792
image observations) while structural supervision is sparse (available for
exactly 48 of those observations, one per specimen).
\item Repeated image observations of the same specimen are correlated, not
independent, and must be treated at the specimen level for evaluation
purposes.
\item No repeated or intermediate structural measurement exists for any
specimen; the dataset supports estimation of a single terminal structural
state per specimen, not a directly observed structural trajectory or
failure time.
\item Row-count and missing-observation conventions differ across the
project's implementation phases and must be stated explicitly whenever a
quantitative result is cited.
\item Campaign identity, steel-mesh configuration, chloride concentration,
and terminal ageing duration co-vary near-perfectly across the observed
specimens, limiting what can be inferred from pooled associations and
giving cross-campaign evaluation the character of extrapolation rather
than ordinary generalisation.
\end{itemize}

These constraints are the principal reason the project's modelling
strategy, described next in Chapter~\ref{ch:framework}, treats
specimen-level grouping, multiple evaluation regimes, and a conservative
distinction between measured, modelled, and derived-screening quantities
as required elements of the methodology rather than optional refinements.
